% 【研读实验室】所属：第1章（由 chapters/revised/01-knowledge.tex \input 引入）
\minitopic{研读实验室：怎样比较Gettier解法}\index{Gettier案例}\index{无假前提}\index{因果理论}\index{知识分析} 一个新条件能排除原始两个案例，只是最低要求。成熟评价还要检验三件事：它是否错误排除普通知识；是否只是换词重述“没有Gettier问题”；是否能统一处理介入运气与环境运气。逐一修补案例容易形成无限清单，比较框架则迫使理论解释自己的范围和机制。

\minitopic{无假前提：简单但范围有限}无假前提条件要求，主体的推理链不能依赖假命题。职位案例中“琼斯会得到职位”是假，因而史密斯的结论被排除。优点是诊断清楚，也保留了演绎确证传递的大部分结构。问题是，并非所有Gettier案例都包含主体接受的假前提。 停止的钟案例中，主体看到表盘恰好显示正确时间。他未必先相信“这只钟正常”再显式推理；即使把背景预设算作前提，许多普通知觉也依赖可能为假的背景假设。假谷仓中，实际所见就是真谷仓，视觉经验和“那里有谷仓”都真，失败来自环境中大量相似假立面。无假前提因而更像处理一类推理性Gettier，而非知识的一般充分条件。 理论还必须说明“基于”什么。一个主体拥有假信念，但目标信念并不依赖它，不能仅因心理系统中存在错误就取消知识。反之，主体未在意识中表述的错误预设可能真正引导推理。基于关系比信念清单更重要。

\minitopic{击败理论：把缺失信息纳入评价}击败理论问：是否存在一个真信息，若主体适当获得，会撤销其确证？“琼斯不会获职”和“这只钟已经停了”显然具有这种作用。条件与实际探究相连，因为我们经常通过寻找潜在击败者改善信念。 但若任意真命题都可作为击败者，理论会过宽。某个遥远星系此刻爆炸是真，却与我看见桌子无关。击败必须通过证据相关性、说明关系或规范可得性连接目标根据。若要求主体能实际取得全部击败信息，又会对信息受限者过苛。 还要处理误导性击败者。一个一贯可靠的人错误告诉你灯光会使红物看绿，你因而暂停颜色判断；后来校准报告证明灯光正常。错误证言虽然不真，却能合理地心理击败原确证；校准又击败该击败者。知识状态随时间和信息层级变化，不能仅用“是否存在一个真反驳”刻画。

\minitopic{因果与可靠性：连接事实和方法}因果理论要求使$p$为真的事实以适当方式造成主体相信$p$。正常知觉最符合此图景：对象造成经验，经验造成信念。Gettier案例中真值事实常在因果链之外，或通过偏离路径进入。理论的吸引力是把“正确连接”具体化，而不是继续列心理条件。 它的难题包括先验知识、否定事实、普遍命题和因果偏离。相信“没有独角兽”似乎不是由一个单一的“无独角兽事实”造成；数学真理也不以通常事件因果作用。若把因果扩大为任何可靠信息关系，理论又逐渐变成可靠主义。 可靠主义评价过程在相关环境中是否倾向产生真信念\parencite{goldman1979}。它自然说明儿童、动物与无反思主体的知识，也能与心理学和工程可靠性结合。过程类型问题却非常尖锐：描述成“视觉”“清晨雾中看远处红灯”还是“这次看见真实红灯”，可靠率完全不同。类型必须在知道结果前由机制、功能与任务稳定划定。

\minitopic{安全性与能力：从不易错到因能力而对}安全性要求，在主体以相近方式相信$p$的近邻情形中，$p$不轻易为假。它直接排除停止的钟与假谷仓，因为微小环境变化便产生同样信念而命题为假。安全并不要求主体实际排除每个怀疑可能，关键是错误世界是否足够接近。 能力理论进一步区分安全的偶然成功与可归功成功。假设一位新手随意按下一台自动验证设备，而设备非常可靠；信念可能安全，却不明显是新手的认知成就。若知识至少是一种成就，需要说明主体、工具和制度各自贡献。把全部功劳只归给个人，会错误排除证言和分布式认知；把设备任何可靠输出都算作个人能力，又使能力条件空洞。 反运气德性方案组合安全与能力，能覆盖多类案例\parencite{pritchard2012}。代价是理论更复杂，也可能面对条件冲突：在假谷仓中，亨利实际视觉表现正常，成功可部分归功能力，却不安全；在高度受控实验中，团队流程让信念安全，单个成员却未理解结果。组合理论要说明知识归属于个人、团队还是系统。

\begin{center}
\begin{tabularx}{\textwidth}{@{}p{2.2cm}XXX@{}}
\toprule
方案 & 核心问题 & 优势 & 主要压力 \\
\midrule
无假前提 & 推理是否依赖假命题 & 处理原始推理案例清楚 & 无假前提的环境运气 \\
击败 & 是否有撤销根据的信息 & 接近实际调查与更新 & 相关性、误导击败与可得范围 \\
因果/可靠 & 事实与方法如何导向信念 & 可自然化并覆盖无反思知识 & 偏离、过程类型、个案错误抵消 \\
安全/能力 & 是否容易错，是否因能力而对 & 统一反运气与认知成就 & 近邻标准与分布式功劳 \\
\bottomrule
\end{tabularx}
\end{center}

\minitopic{最小对比法：一次只改一个变量}思想实验最有分析力的用法，不是让故事越来越奇特，而是制作最小对比。以停止的钟为基准：A中钟正常；B中钟停了，主体恰在正确时刻看；C中钟停了，但主体另以电话报时核验；D中钟正常，却有人给出可信的故障警告。A与B改变过程可靠性，B与C改变信念基础，A与D改变击败信息。若知识判断变化，就能定位理论需要解释的变量。 对比时要保持目标命题一致。“现在是八点”与“钟显示八点”不同；停钟仍可让后者成为知识。还要保持主体实际根据一致。若加入独立核验，不能说同一信念方法突然更安全，而应承认形成基础已经改变。许多争论来自案例叙述未说明主体究竟基于什么。 读者直觉不一致时，可进一步拆分评价：主体是否合理、是否无责、是否可靠、是否知道、是否值得信任。有人判断“他没有做错什么”，却误以为这等于“他知道”。将问题逐项询问，常能发现分歧不在知识判断，而在责任评价渗入。

\minitopic{实验哲学与Gettier直觉}对Gettier案例的问卷研究可以检验判断在不同语言、文化和措辞下是否稳定，但调查设计必须避免把复杂故事变成阅读理解测试。参与者要明确目标命题、真值、根据与偶然因素；问题要区分“知道”“有充分理由”和“只是运气”。若参与者未注意关键细节，其回答不能直接支持概念差异。 即使多数人把某案例判为知识，理论后果也取决于研究目标。描述普通用法的理论必须认真解释；把知识视为规范性认知成就的理论可以说直觉受责任或概率混淆，但要给出独立理由。少数哲学家的稳定判断也不是自动权威，因为专业训练可能增加辨别力，也可能造成理论污染。 最有价值的结果往往不是投票比例，而是变量效应：当真值路径与理由路径更明显脱节时，知识归属是否下降；当主体加入核验时是否恢复；当运气来自环境而非推理时，判断是否变化。变量效应能帮助理论定位机制。

\minitopic{综合案例：医学报告的双重错配}医院系统把甲、乙两人的样本条码互换。甲确有疾病，乙没有；随后报告发送模块又把收件人互换，于是甲收到写着自己姓名且显示阳性的报告。甲基于报告相信自己患病。命题为真，医院总体检测可靠，疾病也因甲的样本在系统中留下因果痕迹，但正确报告由两次错误抵消。 无假前提理论要决定甲是否相信“这是由我的样本正常生成的报告”这一背景命题；击败理论指出条码交换记录；因果理论面对一条既由甲的疾病启动、又发生偏离的链；可靠主义若只看检测仪准确率会漏掉身份匹配过程；安全性指出任一交换不发生，甲的信念便可能为假；能力理论则问甲只是接受制度输出，知识能否归功于他。 若甲的医生在诊断前重新独立采样，第二次阳性使后来信念获得新的基础。不能把后来的验证倒灌回第一次阅读报告。若医院知道系统近期有错配警报却未暂停报告，甲本人可能无责，机构却有严重认识责任。一个案例可以同时包含个体知识失败与制度失职，二者不应相互替代。

\begin{tcolorbox}[breakable,colback=white,colframe=blue!30!black,title={本章研读任务},fonttitle=\bfseries]
自拟一个Gettier案例并制作四个最小版本：原始运气版、移除运气版、加入真击败者版、加入独立核验版。对每版分别填写JTB、无假前提、击败、可靠、安全与能力理论的判定；若某理论无法判定，说明缺失的是哪项细节。
\end{tcolorbox}
